% concatenated from arXiv-5.tex
\documentclass[11pt]{article}

% ---------- Encoding ----------
\usepackage[utf8]{inputenc}
\usepackage[T1]{fontenc}

% ---------- Math ----------
\usepackage{amsmath, amssymb, amsthm, mathtools}

% ---------- Layout ----------
\usepackage[a4paper, margin=1in]{geometry}

% ---------- Hyperlinks ----------
\usepackage[colorlinks=true, linkcolor=blue, citecolor=blue]{hyperref}

% ---------- Theorems ----------
\newtheorem{theorem}{Theorem}[section]
\newtheorem{proposition}[theorem]{Proposition}
\newtheorem{lemma}[theorem]{Lemma}

\theoremstyle{remark}
\newtheorem{remark}[theorem]{Remark}

% ---------- Commands ----------
\newcommand{\R}{\mathbb{R}}
\newcommand{\C}{\mathbb{C}}
\renewcommand{\refname}{Reference} % Para article, report
% \renewcommand{\bibname}{References} % Para book, report
% ---------- Title ----------
\title{
A new equivalence to the Riemann Hypothesis by means of the Salem integral equation
}

\author{
B.~J. Gonz\'{a}lez$^b$ \thanks{Email: bjglez@ull.edu.es} \and
E.~R. Negr\'{i}n$^{a,b}$ \thanks{Email: enegrin@ull.edu.es} \\
\small $^a$Departamento de An\'alisis Matem\'atico, Universidad de La Laguna. Spain \\
\small $^b$Instituto de Matem\'aticas y Aplicaciones (IMAULL), Universidad de La Laguna. Spain
}

\date{}

\begin{document}

\maketitle

% ---------- Abstract ----------
\begin{abstract}
This note presents a new equivalence to the Riemann Hypothesis by means of the Salem integral equation.
\end{abstract}

\bigskip

\noindent \textbf{Keywords:} Riemann Hypothesis; Salem integral equation

\medskip

\noindent \textbf{MSC (2020):} Primary: 11M06; Secondary: 11M36, 44A15

% ============================================================
\section{Introduction}

\hspace{1 cm} The Riemann Hypothesis states that all non-trivial zeros on $0<\operatorname{Re} s<1$ of the Riemann zeta-funcion $\xi (s)$ lie in the line $\operatorname{Re} s=\frac{1}{2}$.

The Rapha\"el Salem publication of 1953 \cite{Salem} proves that the Riemann Hypothesis is true if and only if for each $\frac{1}{2}< \delta <1$ the only bounded measurable function $f(t)$ on $(0,\infty )$ satisfying
\[
\int_{0}^{\infty}\frac{t^{\delta -1}}{e^{xt}+1} f(t) dt =0,\quad \textrm{for\;all} \;x>0 ,
\] 
is $f=0$ almost everywhere.



In this note we prove that the Riemann Hypothesis is true if and only if for each $\frac{1}{2}<\delta <1$  the integral equation
\[
\int_{0}^{\infty}\frac{t^{\delta -1}}{e^{xt}+1} f(t) dt =0,\quad \textrm{for\;all}\;x>0,
\]
has not solutions of type $f(t) = t^{i\gamma}$, $t\in (0,\infty )$, $\gamma \in\mathbb{R}$.

Observe that $f(t) = t^{i\gamma}=e^{i\gamma \, \ln \, t}$ is a bounded measurable function on $(0,\infty )$ whose modulus is equal to one.


% ============================================================
%\section{Main Result}
\section*{Proof of the new equivalence}

\hspace{1 cm} For $\operatorname{Re} s>1$, the Riemann zeta-function $\xi (s)$  is represented by the convergent series:
\[
\xi (s)=\sum_{n=1}^{\infty}\frac{1}{n^s} .
\]

Also, denote by $\eta (s)$ the Dirichlet eta-function represented for $\operatorname{Re} s>0$ by the convergent series:
\[
\eta (s)= \sum_{n=1}^{\infty}\frac{(-1)^{n-1}}{n^s} .
\]

For $\operatorname{Re} s >1$ the next known relation holds
\[
\eta (s)=\left( 1- 2^{1-s}\right) \xi (s) ,
\]
which extends $\xi (s)$ to $0<\operatorname{Re} s<1$.

On the other hand, it is known that for $\operatorname{Re} s >0$
\begin{equation}\label{1.1}
\int_{0}^{\infty}\frac{t^{s -1}}{e^{xt}+1}  dt =x^{-s}\Gamma (s) \eta (s) ,\quad \textrm{for\;all} \: x>0.
\end{equation}


Thus, we obtain that $s= \delta + i\gamma$, $0<\delta < 1$, $\gamma\in\mathbb{R}$, is a zero of $\xi (s) $ if and only if 
\[
\int_{0}^{\infty}\frac{t^{\delta -1}}{e^{xt}+1} t^{i\gamma }  dt = 0,\quad \textrm{for\;all}\: x>0.
\]


Therefore, the Riemann Hypothesis is true if and only if for each $0<\delta < 1$, $\delta \neq \frac{1}{2}$, the integral equation
\[
\int_{0}^{\infty}\frac{t^{\delta -1}}{e^{xt}+1} f(t)  dt = 0, \quad \textrm{for\;all}\: x>0,
\]
has not solutions of type $f(t)=t^{i\gamma}$, $t\in (0,\infty )$, $\gamma\in \mathbb{R}$.


So, from the symmetry of the zeros of $\xi (s)$ on $0<\operatorname{Re} s<1$ with respect to the line $\operatorname{Re}  s =\frac{1}{2}$, one arrives to the next equivalence:

The Riemann Hypothesis is true if and only if for each $\frac{1}{2} < \delta < 1$ the integral equation 
\[
\int_{0}^{\infty}\frac{t^{\delta -1}}{e^{xt}+1} f(t)  dt = 0, \quad \textrm{for\;all}\: x>0,
\]
has not solutions of type $f(t)=t^{i\gamma}$, $t \in (0,\infty )$, $\gamma\in \mathbb{R}$.

\section*{A consequence}

\hspace{1 cm} Observe that from the Salem equivalence \cite{Salem} and from this new equivalence one has:

For each $0< \delta <1$, $\delta \neq \frac{1}{2}$, the only bounded measurable function $f(t)$ on $(0,\infty )$ satisfying
\[
\int_{0}^{\infty}\frac{t^{\delta -1}}{e^{xt}+1} f(t) dt =0,\quad \textrm{for\;all} \;x>0 ,
\] 
is $f =0$ almost everywhere if and only if for each $0<\delta <1$, $\delta \neq \frac{1}{2}$,
the integral equation
\[
\int_{0}^{\infty}\frac{t^{\delta -1}}{e^{xt}+1} f(t) dt =0,\quad \textrm{for\;all}\;x>0,
\]
has not solutions of type $f(t) = t^{i\gamma}$, $t\in (0,\infty )$, $\gamma \in\mathbb{R}$.




% ============================================================
\begin{thebibliography}{99}

\bibitem{Salem}
Rapha\"el Salem,
\textit{Sur une proposition équivalente à l'hypothèse de Riemann},
C. R. Acad. Sci. Paris \textbf{236} (1953), 1127-1128.



\end{thebibliography}

\end{document} 


Let's take $s= \frac{1}{2}+ i\gamma _0$, $\gamma _0\in\mathbb{R}$ that it is a zero of the function $\eta $ for $\Re s >0$, $\Re s \neq 1$ (which we know exists).

Now, from (\ref{1.1}) it follows
\[
\int_{0}^{+\infty}\frac{t^{-\frac{1}{2}+ i\gamma _0}}{e^{xt}+1}  dt = x^{-\frac{1}{2}- i\gamma _0} \Gamma (\frac{1}{2}+ i\gamma _0 )\left( 1 - 2^{-\frac{1}{2}- i\gamma _0}\right) \eta \left( \frac{1}{2}+ i\gamma _0\right) .
\]


We consider the function
\[
f(t)=t^{i\gamma _0} = e^{i\gamma _0 \,\ln t} ,\quad t>0,
\]
whose module is the unit and therefore the function is bounded.

Then, one has
\[
\int_{0}^{+\infty}\frac{t^{-\frac{1}{2}}}{e^{xt}+1} t^{i\gamma _0}dt = 0
\]
for all $x>0$, since $\xi (\frac{1}{2}+i\gamma _0)=0$.

\section*{Main Result: New equivalence to te Riemann Hypothesis}

Taking into account that for $\gamma \in\mathbb{R}$ one has $\delta + i\gamma$ is a zero of the function $\eta$ de Dirichlet if and only if there exists a $\delta >0$, $0<\delta < 1$, such that 
\[
\int_{0}^{+\infty}\frac{t^{\delta -1}}{e^{xt}+1}  t^{- i\gamma }dt =0,
\]
for all $x>0$.

The Riemann Hypothesis is false if there exist a $\delta >0$, $\frac{1}{2} < \delta < 1$ and there exists a $\gamma \in \mathbb{R}$ such that
\[
\int_{0}^{+\infty}\frac{t^{\delta -1}}{e^{xt}+1}  t^{- i\gamma }dt =0,
\]
for all $x>0$.

Note that for $t^{i\gamma} = e^{i\gamma \,\ln t}$ it follows $|t^{i\gamma}| = | e^{i\gamma \,\ln t}| =1$.

From which we concludes that the Riemann Hypothesis is true if and only if for all $\frac{1}{2} <\delta < 1$, the integral equation
\[
\int_{0}^{+\infty}\frac{t^{\delta -1}}{e^{xt}+1} f(t) dt =0,\quad \:\: \textrm{for all} \: x>0 ,
\]
has not solutions of type $f(t)=t^{i\gamma}$, $\gamma\in \mathbb{R}$.

\section{Conclusions}

As a consequence of all the above, it follows that for all $\frac{1}{2} <\delta <1$ the integral equation
\[
\int_{0}^{+\infty}\frac{t^{\delta -1}}{e^{xt}+1} f(t) dt =0,\quad \textrm{for all} \: x>0 ,
\]
has not bounded measurable solutions other than the trivial one if and only if for $\frac{1}{2} <\delta <1$ this integral equation has not solutions of type $f(t) =t^{i\gamma}$, $\gamma\in\mathbb{R}$


\section{A Spectral Operator Associated with the Zeros of $\Lambda(s)$}

In this section, we introduce an operator whose spectrum encodes the zeros of the function
\[
\Lambda(s) := \Gamma(s+1)(1-2^{1-s})\zeta(s).
\]

Let
\[
Z_\Lambda := Z_\zeta \cup Z_p,
\]
where $Z_\zeta$ denotes the nontrivial zeros of $\zeta(s)$ and $Z_p$ those of the prefactor $(1-2^{1-s})$, excluding $s=1$.

\subsection{Definition of the operator}

We consider the logarithmic variable $t = e^u$, mapping functions on $(0,\infty)$ to functions on $\mathbb{R}$. Let
\[
\mathcal{F}: L^2(\mathbb{R}) \to L^2(\mathbb{R})
\]
denote the Fourier transform.

We define a (formal) operator $\widehat{R}$ acting on suitable functions $\phi \in L^2(\mathbb{R})$ by
\[
\widehat{\widehat{R}\phi}(\xi) = m(\xi)\widehat{\phi}(\xi),
\]
where the symbol $m(\xi)$ is defined by
\[
m(\xi) := i\left(\frac{1}{2} - \lambda(\xi)\right),
\]
and $\lambda(\xi)$ ranges over the zeros of $\Lambda(s)$.

\subsection{Spectral structure}

Formally, the spectrum of $\widehat{R}$ is given by
\[
\sigma(\widehat{R}) = \{ i(1/2 - \lambda) : \lambda \in Z_\Lambda \}.
\]

Thus, the zeros of $\Lambda(s)$ determine the spectral values of $\widehat{R}$.

\subsection{Connection with Mellin transform}

Under the Mellin transform
\[
f(t) \longleftrightarrow \phi(u), \quad t=e^u,
\]
functions of the form
\[
f(t)=t^{-1/2-i\gamma}
\]
correspond to plane waves
\[
\phi(u)=e^{i\gamma u}.
\]

These functions can be viewed as generalized eigenfunctions associated with the spectral values
\[
i(1/2 - (1/2+i\gamma)) = \gamma.
\]

\subsection{Interpretation}

The operator $\widehat{R}$ provides a spectral model in which:
\begin{itemize}
\item the zeros of $\Lambda(s)$ correspond to spectral parameters,
\item the critical exponent $1/2$ appears as a natural shift,
\item the Mellin transform diagonalizes the associated integral equation.
\end{itemize}

\subsection{Relation with the main result}

The triviality results established in this work can be interpreted in terms of the operator $\widehat{R}$.

\begin{proposition}
Let $f$ satisfy the assumptions of the main theorem. Then its Mellin transform has no support on the spectral set associated with $\widehat{R}$.
\end{proposition}

\begin{proof}
The growth condition
\[
f(t)=O(t^{1/2+\varepsilon})
\]
excludes contributions of the form $t^{-1/2-i\gamma}$, which correspond to generalized eigenfunctions of $\widehat{R}$. Therefore, no spectral component associated with $\sigma(\widehat{R})$ can be present.
\end{proof}

\begin{remark}
The operator $\widehat{R}$ is introduced as a spectral model and is not constructed here as a self-adjoint operator on a specified dense domain. A fully rigorous realization would require additional functional-analytic developments.
\end{remark}
% ============================================================
%\section*{Acknowledgements}

% Optional
